\subsubsection{An integral form with unit constant}
\label{par:count-three-cusp-correction}

Divisibility of one Eisenstein constant is only the first step toward a cusp form. An integral correcting form must change that constant while prescribing its values independently at the other cusps.

\begin{proposition}[Cusp correction with unit constant term]\label{prop:unit-constant-cusp-correction}
Let \(N\ge3\), \(\ell\) be odd and prime to \(N\), and
\(\chi(-1)=(-1)^k\). With the separating weight \(m\) of
\eqref{eq:count-three-correction-degree}, assume \(k\ge m+10\).
For the correcting form \(F\), let \(a_n(F;c)\) denote the expansion coefficients at a cusp \(c\), in its chosen integral coordinates. Then 
\[
 a_n(F;c)\in O\quad\text{at every cusp},\qquad a_0(F;\infty)=1,
\]
 with \(O\) the coefficient DVR. This is
the correction used for the primitive Eisenstein series.
\end{proposition}

\begin{proof}
Let \(N\geq3\) with \(\ell\nmid N\), and let
\(\zeta\) be a primitive \(N\)-th root of unity. We use its complex
value for Fourier expansions and its \(\ell\)-adic value for
integrality. For \(a\in(\Z/N)^\times\) and \(r\geq3\), take the
normalized weight-\(r\) Eisenstein sum over \((c,d)\equiv(0,a)\bmod N\)
with summand \((cz+d)^{-r}\). Denote it by \(E_{r,a}\). The normalization is
\[
 E_{r,a}(z)=\frac{N^r(r-1)!}{(-2\pi i)^r}
          \sum_{(c,d)\equiv(0,a)\bmod N}(cz+d)^{-r}.
\]
Absolute convergence gives holomorphy and the right slash action of
\(\gamma=\bigl(\begin{smallmatrix}*&*\\c_\gamma&d_\gamma\end{smallmatrix}\bigr)\)
transports the row \((0,a)\) to \((ac_\gamma,ad_\gamma)\).

In particular the form is modular on \(\Gamma_1(N)\). For
\(\gamma\in\Gamma_0(N)\), the slash is \(E_{r,ad_\gamma}\).
If \(c_\gamma\not\equiv0\bmod N\), the transported sum contains
only terms with \(c\ne0\). Summing the absolute values over \(d\)
gives a bound \(O_r((|c|\operatorname{Im} z)^{1-r})\) for
\(\operatorname{Im} z\geq1\). Writing \(y=\operatorname{Im}z\ge1\) gives the explicit tail bound
\[
 \sum_{c\ne0}(|c|y)^{1-r}
 =2y^{1-r}\sum_{c\ge1}c^{1-r}\longrightarrow0
 \qquad(y\to\infty),
\]
since \(r>2\). The same bound on the
nonconstant part and the convergent \(c=0\) sum give holomorphy
at every cusp.

Set \(L_r(x)=(x\,d/dx)^{r-1}(x/(1-x))\). Differentiating the
partial fraction expansion of \(\pi\cot(\pi t)\) gives
\[
  \frac{(r-1)!}{(-2\pi i)^r}
      \sum_{b\in\Z}(t+b)^{-r}
     =L_r(e^{2\pi i t}).
\]
For \(\operatorname{Im}t>0\) the right side has expansion
\(\sum_{e\geq1}e^{r-1}e^{2\pi i et}\). For the constant
term take \(t=a/N\). For \(c>0\) use \(t=(cz+a)/N\).
The terms with \(c<0\) give the contribution with \(-a\)
and the factor \((-1)^r\). Thus at infinity
\[
 a_0(E_{r,a})=L_r(\zeta^a),\qquad
 a_n(E_{r,a})=\sum_{e\mid n}e^{r-1}
              \bigl(\zeta^{ae}+(-1)^r\zeta^{-ae}\bigr)
       \quad(n\geq1).
\]
The numerator of \(L_r\) is integral and its denominator is
\((1-x)^r\). Reduction is injective on the \(N\)-th roots of
unity, so all these coefficients are integral at \(\ell\).

We next separate the constant at \(a=1\) from that at any
\(a=d\ne\pm1\), with the same weight in the two evaluations.
In the residue field put \(t=x/(1-x)^2\) and
\(u=(1+x)/(1-x)\) for \(x=\bar\zeta^a\). Direct differentiation gives
\[
 L_3(x)=tu,\qquad L_4(x)=t(1+6t),\qquad
 L_5(x)=tu(1+12t),\qquad u^2=1+4t.
\]
At any unit label \(L_3(x)\ne0\), because \(x\ne1,-1\).

In residue characteristic other than two or three, equality of
both \(L_3^2\) and \(L_3L_5\) at \(x,y\) forces
\(t(x)=t(y)\) by taking their quotient.
In characteristic three \(L_4=t\) suffices. Clearing denominators
in \(t(x)=t(y)\) gives \((x-y)(1-xy)=0\).
For \(x=\bar\zeta\), \(y=\bar\zeta^d\), both alternatives
are excluded. At least one of the three families
\[
      E_{4,a},\qquad E_{3,a}^2,\qquad E_{3,a}E_{5,a}
\]
therefore has constants at \(a=1,d\) differing by a unit.

For each such \(d\) choose one family \(H_{d,a}\) of weight \(r_d\).
For a level-one form \(L_{r_d}\) of weight \(r_d\)
with integral coefficients and constant one, set
\[
        H'_{d,a}=H_{d,a}-a_0(H_{d,d})L_{r_d}.
\] Such
forms exist for all even nonnegative weights except two,
by products of \(E_4\) and \(E_6\), including the empty product.
The resulting family has zero constant at \(a=d\) and unit
constant at \(a=1\). Multiplication of the correcting factors gives the even weight
\begin{equation}\label{eq:count-three-correction-degree}
                 m=\sum_{d\ne\pm1}r_d\le8\varphi(N).
\end{equation}
It has zero constants at all \(a\ne\pm1\) and retains a
unit at \(a=1\). The same right slash rule holds for all the
families, since the coefficients in these operations are fixed
as \(a\) varies.

For \(k\geq m+10\), multiply this family by \(E_{3,a}\) if
\(k\) is odd and by \(E_{3,a}^2\) if \(k\) is even. Complete the weight with a product of \(E_4\) and \(E_6\) having constant term one.
The forms \(F_a\) of weight \(k\) so obtained vanish at all
cusps outside the \(\Gamma_0(N)\)-orbit of infinity by the
weight-three factor. At infinity \(c_1=a_0(F_1)\) is a unit. The constant terms
are
\[
       a_0(F_a)=
       \begin{cases}
       c_1,&a=1,\\
       (-1)^kc_1,&a=-1,\\
       0,&a\ne\pm1 .
       \end{cases}
\]
The second case is obtained by slashing by \(-I\).

For \(\chi(-1)=(-1)^k\), set
\[
    G=\frac{1}{2c_1}\sum_{a\in(\Z/N)^\times}\chi(a)^{-1}F_a .
\]
The right slash rule and reindexing \(a\mapsto ad_\gamma\)
give \(G|_k\gamma=\chi(d_\gamma)G\) on \(\Gamma_0(N)\).
Its coefficients are integral because \(2c_1\) is a unit.
Its constant is one, and it vanishes away from the orbit of
infinity. In particular the construction works at the level
\(N=Mw\) with \(w\equiv1\bmod\ell\), even though
\(\ell\mid\varphi(N)\).
\label{cor:ordinary-count-iii-eigenform-1}

\end{proof}

\subsubsection{The cuspidal congruence}

The integral correction can now be multiplied by the small Eisenstein constant. Subtracting it enforces cuspidality while preserving the positive Fourier coefficients modulo the maximal ideal.

\begin{proposition}[A cuspidal congruence at the chosen conductor]\label{prop:count-three-cusp-congruence}
For the primitive Eisenstein character \(\chi\) of conductor \(N=Mw\)
and the small constant calculated for the seed, the cusp correction gives an integral cusp form \(f\) satisfying 
\[
 a_n(f)\equiv\sum_{e\mid n}\chi(e)e^{k-1}\pmod{\mathfrak m}\qquad(n\ge1).
\]
 The
congruence includes the ordinary Hecke eigenvalue at \(\ell\)
used in the lift.
\end{proposition}

\begin{proof}
For the primitive character \(\chi=\phi\theta\) of conductor
\(N=Mw\), the weight-\(k\) Eisenstein series \(E\) of character
\(\chi\) is the sum of \(E_{k,a}\) with weights
\(\chi(a)^{-1}/(2\mathcal G(\chi^{-1}))\) over the unit labels.
Here \(\mathcal G(\chi^{-1})=\sum_a\chi(a)^{-1}\zeta^a\)
is the nonzero primitive Gauss sum, with characters extended
by zero on nonunits. Primitivity gives
\(\sum_a\chi(a)^{-1}\zeta^{ae}=\chi(e)\mathcal G(\chi^{-1})\)
also when \(e\) is a nonunit.

For a nonunit, multiplying
\(a\) by an element fixing \(e\) and on which \(\chi\) is
nontrivial gives zero by primitivity. Parity then gives
\[
    a_n(E)=\sigma(n):=\sum_{e\mid n}\chi(e)e^{k-1}
    \quad(n\geq1),\qquad a_0(E)=L(1-k,\chi)/2 .
\]
For the constant, the same Gauss-sum identity applied
to \(\sum_{e\geq1} e^{k-1}\zeta^{ae}x^e\) for \(|x|<1\)
gives the finite sum of \(L_k(\zeta^a x)\).

At \(x=e^t\), the generalized Bernoulli generating
function gives the continued value
\[
 \left.\left(\frac{d}{dt}\right)^{k-1}
            \frac{\sum_{e=1}^N\chi(e)e^{et}}{1-e^{Nt}}
       \right|_{t=0}=L(1-k,\chi).
\]
The primitive Gauss factor and parity give the normalization of
the constant term above. The cusp vanishing and the character follow
from the corresponding transformations of \(E_{k,a}\).
Set \(f=E-a_0(E)G\). Its constant term at infinity is zero,
and the other positive coefficients satisfy
\[
      a_n(f)\equiv a_n(E)=\sigma(n)\pmod{\mathfrak m}\qquad(n\ge1)
\]
because \(a_0(E)\in\mathfrak m\) and \(G\) has integral coefficients. At cusps in the orbit of infinity,
the vanishing follows by applying the common character law
to the zero constant of \(f\) at infinity. Here the smallness
of \(a_0(E)\) is the earlier Bernoulli and Euler-factor
calculation for the chosen \(k\).

For a prime \(p\), the Hecke operator in this character space has
coefficient \(a_{pn}+\chi(p)p^{k-1}a_{n/p}\), where the second
term is zero if \(p\nmid n\), and \(\chi(p)=0\) if \(p\mid N\).
The divisor sums satisfy the same formula with eigenvalue
\(1+\chi(p)p^{k-1}\). In fact, if \(n=p^b n'\) with
\(p\nmid n'\) and \(v=\chi(p)p^{k-1}\), this is the polynomial
identity
\[
  (1+v+\cdots+v^{b+1})+v(1+v+\cdots+v^{b-1})
      =(1+v)(1+v+\cdots+v^b),
\]
with the second sum empty at \(b=0\).
Thus the reduction of \(f\) is a simultaneous Hecke
eigenvector, nonzero since \(a_1(f)\equiv1\).

\end{proof}

\subsubsection{Simultaneous lifting and the ordinary line}

A nonzero mod-\(\ell\) cusp form gives the desired residual eigenvalues. The second tensor factor requires one characteristic-zero eigenform realizing them simultaneously, including the ordinary eigenvalue at \(\ell\).

\begin{proposition}[Eigensystem of the corrected cusp form]\label{prop:count-three-cuspidal-eigensystem}
The cuspidal congruence of Proposition~\ref{prop:count-three-cusp-congruence}
gives the simultaneous cuspidal eigenform \(g\) and the ordinary local line used for the rank-two factor. After a finite coefficient extension the normalized eigenform satisfies 
\[
 a_p(g)\equiv1+\chi(p)p^{k-1}\pmod{\varpi}\qquad(p\text{ prime}).
\]
 The weight can be
fixed before the choice of the final deep layer.
\end{proposition}

\begin{proof}
\proofstep{An ordinary eigensystem with the selected residual line.}

Choose a finite local coefficient field \(E_\ell\) containing the
coefficient fields of the forms in the construction, and let \(O\)
be its valuation ring with uniformizer \(\varpi\). The
character-\(\chi\) cusp space of level \(N\) has a model over
this field. Indeed normalized newforms have algebraic
coefficients in number fields, and the newform decomposition
spans the cusp space. Since the character has conductor \(N\),
every constituent is new at level \(N\). Both the character
and the coefficients of \(f\) are in the chosen coefficient
field, so \(f\) belongs to this model. One can check this by
taking finitely many Fourier coefficients that identify the
coordinates in a newform basis.

Let \(V_O\) consist of the forms with all coefficients in \(O\).
Take a finite list \(n_1,\ldots,n_s\) separating the forms. Its coefficient map gives
\[
 V_O\hookrightarrow O^s,\qquad h\longmapsto
 (a_{n_1}(h),\ldots,a_{n_s}(h)).
\]
A submodule of a finite free module over a DVR is finite free.
It also contains a basis after integral scaling, by
integrality of the normalized newforms. It is finite free
over \(O\). The coefficient formula makes it stable under
all prime Hecke operators, including those at \(p\mid N\). The coefficient condition gives the saturated reduction
\[
       \ker(V_O\longrightarrow \text{residue expansions})
                                      =\varpi V_O.
\] The mod-\(\varpi\) eigenline of \(f\) therefore
defines a character of the commutative \(O\)-algebra of
Hecke operators on \(V_O\). Write \(\mathfrak m\) for its
kernel. This algebra is finite and torsion-free over \(O\).
Write \(\mathbb T\) for the Hecke algebra. A minimal
prime \(\mathfrak p\subset\mathfrak m\) avoids the uniformizer
since it is a nonzerodivisor.

Normalize \(\mathbb T/\mathfrak p\)
in its fraction field and choose a place over \(\mathfrak m\);
this gives one characteristic-zero evaluation
\[
 \mathbb T\longrightarrow
      \operatorname{Frac}(\mathbb T/\mathfrak p)
\]
with all the prescribed residual Hecke values.
In the generic fibre the faithful module has a nonzero
factor at \(\mathfrak p\). After extending scalars, the
maximal ideal of that Artinian local factor acts nilpotently.
If \(M_\mathfrak p\) is this local generic factor and
\(\mathfrak n\) its maximal ideal, choose the last nonzero
\(\mathfrak n^a M_\mathfrak p\). It is killed by \(\mathfrak n\),
so a nonzero vector there is a common eigenvector.
The eigenvector has nonzero first coefficient, since the
Hecke recursion from a zero first coefficient would give
zero coefficients at every positive index. Normalizing it
gives a cuspidal eigenform \(g\) of level \(N\) and character
\(\chi\) with
\[
    a_p(g)\equiv1+\chi(p)p^{k-1}\pmod{\varpi}
       \qquad(p\text{ prime}),
\]
where \(\varpi\) now denotes a uniformizer after coefficient
extension. This is the simultaneous lifting argument of
Deligne and Serre. \cite[Lemma 6.11]{deligne-serre}

Use the Galois representation of the weight-\(k\) newform
\(g\), with arithmetic Frobenius polynomial
\(X^2-a_p(g)X+\chi(p)p^{k-1}\) for \(p\nmid\ell N\).
It is absolutely irreducible, with determinant
\(D=\Phi\theta\). Since \(\theta\) has order \(\ell\), its residual character is trivial. At every good prime the trace congruence and determinant now give
\[
 \bar P_p(X)=X^2-(1+\bar\Phi(\operatorname{Frob}_p))X+
 \bar\Phi(\operatorname{Frob}_p)
 =(X-1)(X-\bar\Phi(\operatorname{Frob}_p)).
\]
Chebotarev and semisimplification yield \(\bar\Phi\oplus1\). For \(p\mid N\) the local component has conductor
exponent equal to that of the central character, by primitivity
of \(\chi\).

It is principal series with one unramified factor.
Indeed, for a supercuspidal parameter the two-dimensional
irreducible Weil representation has invariants of dimension
zero or two on each normal ramification group. Its determinant
has conductor at most half as large by the Artin conductor
formula. For a special representation the conductor is one
with unramified central character if the twist is unramified.
For a ramified twist of conductor exponent \(a>0\) the
conductor is \(2a\), while that of its central character is
at most \(a\). For two principal series characters with conductors
\(a_1,a_2\), the conductor and determinant bounds read
\[
                a(W)=a_1+a_2,\qquad a(\det W)\le\max(a_1,a_2).
\]
At a ramified prime equality thus forces precisely one
\(a_i\) to be zero. Local--global
compatibility gives the two inertia characters \(1,\phi\theta\).
Outside \(\ell N\) the representation is unramified.
\cite{deligne-galois,carayol-local}

Since \(k\geq3\) and \(\ell\nmid N\), the congruence at
\(p=\ell\) gives \(a_\ell(g)\equiv1\). The unit root
\(\alpha_\ell\) of the Hecke polynomial at \(\ell\) is
congruent to one. The ordinary local theorem for this
newform, or for its unit-root \(\ell\)-stabilization, gives
an unramified quotient \(\nu\) with
\(\nu(\operatorname{Frob}_\ell)=\alpha_\ell\) and invariant
line of character \(D\nu^{-1}\). Its inertia character is
\(\epsilon_\ell^{k-1}\). The two generators of its local data satisfy
\[
         \bar\nu|_{I_\ell}=1,\qquad
         \bar\nu(\operatorname{Frob}_\ell)=1 .
\]
Hence \(\bar\nu\) is trivial on the decomposition group. The high-weight line thus reduces to \(\bar\Phi\)
on that entire group, including when
\(k-1\equiv0\bmod\ell-1\).
For the dual of the attached representation the
zero-exponent line has residual character one and determinant
\(D^{-1}\); the twist by \(D\) gives exactly the preceding
order of the characters. \cite{hida-ordinary-galois}\label{cor:ordinary-count-iii-eigenform-2}

The weight was chosen after \(N\) and the initial positive
exponent \(m_0\) of the quadratic factor. In addition to
the cusp-correction cutoff \(m+10\), we may impose
\(k\geq18m_0+3\) and the earlier congruence and coprimality
conditions. No bound on the eventual deep layer is used in
this choice.
\end{proof}
